% =============================================================================
% Abstract for IEEE Paper: Bat-Inspired Drone Navigation
% Hebrew primary, English technical terms in \en{}
% XeLaTeX : requires \selectlanguage{hebrew} at top
% =============================================================================

\selectlanguage{hebrew}


מאמר זה מציג מערכת ניווט אוטונומית לרחפנים בהשראת עטלפים, המשלבת מיזוג חיישנים רב-מודאלי, \en{SLAM} אקוסטי, תכנון מסלול אדפטיבי, ולמידה עמוקה לעיבוד הדים. המערכת מבוססת על מודל ביומימטי של מערכת האקולוקציה העטלפית, הכולל פיצוי תדר דופלר אדפטיבי, סריקת ראש, ומודל קרינת כנפיים. מיזוג החיישנים מבוסס \en{EKF} עם אומדן קווריאנס אדפטיבי מאפשר זיהוי השפלת חיישן באמצעות מבחן כי-ריבוע והתאמת משקל החיישן באומדן המצב בהתאם. אלגוריתם \en{SLAM} אקוסטי בונה מפת רשת תאים הסתברותית תלת-ממדית תוך שימוש ברשת נוירונים גרף (\en{GNN}) להתאמת הדים. מתכנן מסלול מבוסס \en{RRT*} מותאם לאי-ודאות חישה משתנה ומשלב פונקציית עלות רב-מודאלית הכוללת בטיחות, יעילות אנרגטית, ואיכות חישה. תוצאות ניסויים בסביבות מאתגרות (עשן, חושך, עלווה) הראו כי המערכת משיגה \en{RMSE} מיקום של 0.18-0.28 מטר, שיפור של 34-49\% לעומת שימוש בחיישן על-קולי בלבד. שיעור ההימנעות מהתנגשות היה 82-94\% תלוי בסביבה. המערכת הוטמעה על פלטפורמת \en{NVIDIA Jetson Orin NX} ופועלת בזמן אמת עם צריכת מעבד של 42\% במצב ביצועים גבוהים. המאמר מדגים כי השראה ביולוגית ממערכת האקולוקציה של עטלפים יכולה להוביל לפתרונות ניווט חדשניים לרחפנים בסביבות מוגבלות ראות, שבהן מערכות קונבנציונליות מבוססות ראייה או \en{LiDAR} נכשלות.


% =============================================================================
% End of Abstract
% =============================================================================